\documentclass[12pt]{article}
\usepackage[utf8]{inputenc}
\usepackage[T1]{fontenc}
\usepackage{amsmath,amssymb,tikz}
\usepackage{geometry}
\geometry{margin=2.5cm}

\begin{document}

\noindent\textit{[fragment neclar la începutul paginii: $(\mathbb{N},+,\cdot,\leq,1)$ -- relația cu restul paginii nu e clară; posibil rest dintr-o idee neterminată de pe pagina anterioară.]}

\bigskip
\noindent\underline{Lemă}. Fie $(G,\cdot,e)$=grup, $a\in G$, $n\in\mathbb{N}^*$

\noindent U.A.S.E:

\begin{enumerate}
\item[(1)] ord $a=n$.
\item[(2)] (i) $a^n=e$,

\noindent (ii) $a^h=e$, $h\in\mathbb{Z} \Rightarrow n\mid h$
\end{enumerate}

\bigskip
\noindent (1)$\Rightarrow$(2). \ evident.

\bigskip
\noindent (2)$\Rightarrow$(1) \quad $h=qn+r$ \quad $0\leq r<n$.

\[
a^h = a^{qn+r} = a^{qn}\cdot a^r = (a^n)^q\cdot a^r = a^r
\]
\[
a^h=e
\]
\[
\Rightarrow a^r=e
\]

\noindent\footnote{în original, o linie de forma ,,dacă $r>0$, $a^r=e\Rightarrow r=0$'' apare tăiată/eliminată aici, înlocuită cu argumentul de mai jos.}

\bigskip
\begin{center}
$\big(a^n=e$ şi $0<t<n \Rightarrow a^t\neq e\big)$
\end{center}

\noindent pp.\ $a^t=e=a^n \Rightarrow a^{n-t}=e \overset{\text{Lemă}}{\Longrightarrow} n\mid n-t$ \ imposibil.

\begin{center}
\begin{tikzpicture}[scale=1]
\draw (0,0) -- (4,0);
\draw (0,-0.1) -- (0,0.1) node[above] {};
\draw (0,0) node[below] {$0$};
\draw (1.5,-0.15) -- (1.5,0.15);
\draw (1.5,0) node[below] {$t$};
\draw (3,-0.15) -- (3,0.15);
\draw (3,0) node[below] {$n$};
\draw (1.5,0.15) -- (3,0.15);
\end{tikzpicture}
\end{center}

\bigskip
\noindent\underline{Prop}: Fie $(G,\cdot)$=grup, $a\in G$, ord$(a)=n$.

\noindent $h\in\mathbb{Z}$. Atunci ord $a^h = \dfrac{n}{(n,h)}$.

\bigskip
\noindent În particular ord $a^h=n \Leftrightarrow (n,h)=1$.

\bigskip
\noindent\underline{Dem}: Fie $d=(h,n)$ = c.m.m.d.c.

\[
h=dh_1, \quad n=dn_1
\]

\noindent\underline{Lema}: ord $a^h = n_1$ \qquad $b=a^h$.

\[
b^{n_1} = a^{hn_1} = a^{dh_1n_1} = a^{nh_1} = (a^n)^{h_1} = e.
\]

\noindent pp.\ $b^t=e=a^{ht} \Rightarrow n\mid ht \Rightarrow n_1d\mid dh_1t \Rightarrow n_1\mid h_1t$

\noindent dar $(h_1,n_1)=1 \Rightarrow n_1\mid t$

\[
\Rightarrow \text{ord } b=n_1 \Leftrightarrow \text{ord } a^h = \frac{n}{(n,h)}
\]

\end{document}
